\documentclass[11pt]{article}
\usepackage[utf8]{inputenc}
\usepackage[margin=1in]{geometry}

\title{A Gravity-Based Framework for Scaling Emissions Assessment of Drone and E-bike Food Delivery Systems}
\author{Jingxian Lin\\Department of Civil Engineering\\The University of Hong Kong}
\date{\today}

\begin{document}

\maketitle

\begin{abstract}
  China's booming digital economy has driven explosive growth in the food delivery sector, making its environmental impact an increasingly critical concern. As a potentially sustainable alternative, drone delivery offers speed and flexibility and has attracted considerable research interest. However, most existing studies estimating emissions from e-bike or drone delivery lack real origin-destination (OD) data and rely predominantly on aggregate statistical metrics, which fail to capture route-level emission variations.

  This study addresses these gaps by developing a gravity-based OD generation framework that synthesizes door-to-door delivery demand from merchant monthly sales data and building footprint data, enabling spatially explicit, trip-level emission calculations at an unprecedented scale. Applying this framework nationwide, we present the comprehensive comparison of carbon emissions between e-bike and drone delivery systems across multiple Chinese cities with diverse urban features. Our results reveal that while e-bikes generally produce lower emissions under typical urban conditions, drones can achieve emission reductions in specific scenarios characterized by long delivery distances and high detour ratios, suggesting the feasibility of hybrid drone-e-bike delivery systems that reduce overall emissions. These findings demonstrate that the environmental advantage of drone delivery is highly context-dependent and establish a replicable, scalable methodology for designing environmentally sustainable food delivery systems and assessing emerging delivery technologies at the national scale.
\end{abstract}

\end{document}
